% Part: normal-modal-logic
% Chapter: axioms-systems
% Section: introduction

\documentclass[../../../include/open-logic-section]{subfiles}

\begin{document}

\olfileid{nml}{axs}{int}

\olsection{جان پہچان}

ساڈے کول بنیادی امکانی زبان دی معنیات امکانی نمونیاں راہیں اے۔ ساڈے کول !!a{formula} دے معتبر ہون دا خیال وی اے، یعنی سارے نمونیاں دے سارے جہاناں وچ سچا ہونا، یا نمونیاں یا فریماں دی کسے جماعت دے لحاظ نال معتبر ہونا، یعنی جماعت دے سارے نمونیاں دے سارے جہاناں وچ، یا فریم اُتے ٹکے نمونیاں دے سارے جہاناں وچ، سچا ہونا۔ منطق عام طور تے اعتبار دیاں ایہو جیہیاں معنیاتی پچھاناں نوں !!{derivability} دے ثبوتی خیال نال جوڑدی اے۔ مقصد ایہ اے کہ کسے نظام وچ !!{derivability} ایویں متعین کریے کہ !!a{formula} اُدوں تے صرف اُدوں !!{derivable} ہووے جدوں اوہ معتبر ہووے۔

سبھ توں سادے تے تاریخی طور تے سبھ توں پرانے !!{derivation} دے نظام اوہ نیں جیہناں نوں ہلبرٹ طرز دے یا مسلّماتی !!{derivation} دے نظام آکھدے آں۔ کئی امکانی منطقاں لئی ہلبرٹ طرز دے !!{derivation} دے نظام بناؤنے مقابلتاً سوکھے نیں۔ نظریاتی مطالعے وچ اینہاں نال کم کرنا سوکھا اے، جیویں اینہاں دی صحت مندی تے تمامیت ثابت کرنا۔ پر !!{formula}s ثابت کرن لئی اینہاں نوں ورتنا، مثلاً فطری استنباط دے نظاماں دی ورتوں نالوں، کِتے اوکھا اے۔

ہلبرٹ طرز دے !!{derivation} دے نظاماں وچ کسے !!a{derivation} دا، اک !!a{formula} لئی، مطلب !!{formula}s دی اک لڑی اے جیہڑی کجھ مسلّمیاں توں تھوڑے جیہے استنباطی قاعدیاں راہیں متعلقہ !!{formula} تک پہنچدی اے۔ کیوں جے اسیں چاہندے آں کہ !!{derivation} دا نظام معنیات نال میل کھاوے، ایس لئی سانوں ایہ پکا کرنا پئے گا کہ !!{derivable} فارمولیاں دے سیٹ دا ہر فارمولا ہر نمونے وچ سچا ہووے (یا ہر اوس نمونے وچ سچے ہون جیہدے وچ سارے مسلّمے سچے نیں)۔ پہلاں اسیں امکانی منطقاں دیاں اوہ خاصیتاں کڈھ کے ویکھاں گے جیہڑیاں ایس کم لئی ضروری نیں، یعنی «نارمل» امکانی منطقاں دیاں۔ نارمل امکانی منطقاں لئی صرف دو استنباطی قاعدے منّنے پیندے نیں، قاعدۂ اثباتِ مقدم تے لازمیت دا قاعدہ۔ مسلّمیاں وجوں ہمیشہ صادق فارمولیاں دیاں ساریاں (بدل توں بنیاں) مثالاں، تے متعلقہ امکانی منطق دے حساب نال کجھ امکانی مسلّمے، لیندے آں۔ جے ساڈی دلچسپی صرف سارے نمونیاں دی جماعت وچ اے، تاں وی~$\Ax{K}$\iftag{notprvDiamond}{}{ تے~$\Ax{Dual}$} دیاں بدل نال بنیاں ساریاں مثالاں مسلّمیاں وچ گنّیاں جان گیاں۔ ایہناں نال ای گھٹ توں گھٹ نارمل امکانی منطق~$\Log K$ بندی اے۔

\begin{defn}
\emph{قاعدۂ اثباتِ مقدم} استنباط دا ایہ سانچہ اے۔
\begin{prooftree}
\AxiomC{$!A$}
\AxiomC{$!A \lif !B$}
\RightLabel{\MP}
\BinaryInfC{$!B$}
\end{prooftree}
اسیں آکھدے آں کہ !!a{formula}~$!B$ \emph{نکلدا اے}~!!{formula}s $!A$، $!C$ توں قاعدۂ اثباتِ مقدم نال اُدوں تے صرف اُدوں جدوں $!C \ident !A \lif !B$۔
\end{defn}

\begin{defn}
\emph{لازمیت دا قاعدہ} استنباط دا ایہ سانچہ اے۔
\begin{prooftree}
\AxiomC{$!A$}
\RightLabel{\Nec}
\UnaryInfC{$\Box !A$}
\end{prooftree}
اسیں آکھدے آں کہ !!{formula}~$!B$، !!{formula}s $!A$ توں لازمیت دے قاعدے نال اُدوں تے صرف اُدوں نکلدا اے جدوں $!B \ident \Box !A$۔
\end{defn}

\begin{defn}
مسلّمیاں دے سیٹ~$\Sigma$ توں \emph{!!{derivation}} اوہ !!{formula}s دی لڑی اے جیہڑی $!B_1$، $!B_2$، \dots، $!B_n$ نال لکھی جاندی اے، جیہدے وچ ہر $!B_i$ لئی ہِیٹھ لکھیاں وچوں کوئی اک گل سچی اے۔
\begin{enumerate}
\item اوہ کسے ہمیشہ صادق فارمولے توں بدل نال بنی مثال اے؛ یا
\item اوہ~$\Sigma$ دے کسے !!a{formula} توں بدل نال بنی مثال اے؛ یا
\item اوہ دو !!{formula}s $!B_j$، $!B_k$ توں، جتھے $j$، $k < i$، قاعدۂ اثباتِ مقدم نال نکلدا اے؛ یا
\item اوہ کسے !!a{formula}~$!B_j$ توں، جتھے $j < i$، لازمیت دے قاعدے نال نکلدا اے۔
\end{enumerate}
جے کوئی ایسا !!a{derivation} ہووے جیہدے لئی $!B_n \ident !A$، تاں اسیں آکھدے آں کہ
$!A$، \emph{!!{derivable} اے، $\Sigma$ توں}، علامتاں وچ $\Sigma \Proves !A$۔
\end{defn}

ایس تعریف نال پتا لگے گا کہ !!{derivable} فارمولیاں دا سیٹ اک نارمل امکانی منطق بناؤندا اے، تے ہر !!{derivable} !!{formula} ہر اوس نمونے وچ سچا اے جیہدے وچ ہر مسلّمہ سچا اے۔ !!{derivation}s دی ایس خاصیت نوں \emph{صحت مندی} آکھدے آں۔ اُلٹے رخ نوں \emph{تمامیت} آکھدے آں، تے اینوں ثابت کرنا اوکھا اے۔

\end{document}
